this kernel is contained in that of the composite homomorphism:
@@M0@@
the last homomorphism being defined in degree @@M1@@ by the inclusion of @@M2@@ in @@M3@@. Now the kernel of the composite homomorphism is the principal ideal generated by @@M4@@, which, for reasons of integrality, contains no nonzero homogeneous elements. Thus the first homomorphism is injective, completing the proof in this case.

Return to the general case where @@M5@@ is arbitrary; we may still suppose @@M6@@ free over @@M7@@. Attach the index @@M8@@ to the special case just treated; then, by part i) of the proposition, the situation under consideration is obtained by the base change @@M9@@ from the situation studied above, the homomorphism @@M10@@ being defined by taking a basis of @@M11@@, its image @@M12@@ in @@M13@@, and sending @@M14@@ to @@M15@@. Put @@M16@@, @@M17@@; then @@M18@@; the Koszul complex defined by @@M19@@ is deduced by base change from the Koszul complex defined by @@M20@@; since the latter is a projective resolution of @@M21@@ (working on a suitable open subset), the homology groups of @@M22@@ on this open subset are the @@M23@@. These are the degree @@M24@@ components of the @@M25@@; the @@M26@@ vanish for @@M27@@ by the relations 1.2, so @@M28@@ is a resolution of @@M29@@, which shows that @@M30@@ is regular.

iii) Since the property is local on @@M31@@, we may again suppose @@M32@@ affine and @@M33@@ free of finite type. Resume the notation of ii).

First consider the standard case of ii). The ring @@M34@@ is then noetherian and of finite cohomological dimension, and the algebra @@M35@@
